\subsection{Lexikální analýza}
\label{sec:lexer}

Lexikální analyzátor převádí vstupní text na sekvenci tokenů, které jsou
následně zpracovány parserem. V~kompilátoru \\Roth{} je lexer implementován
v~modulu \texttt{src/lexer.rs} jako struktura \texttt{Lexer}, která si uchovává:
(1) celý vstup jako řetězec, (2) aktuální pozici ve vstupu a~(3) informace o~řádku
a~sloupci pro diagnostiku chyb.

Každý token nese kromě svého typu také původní text (pole \texttt{raw}) a~pozici
\texttt{Position} (řádek, sloupec a~offset ve vstupu). Tato volba zjednodušuje
generování chybových hlášení ve všech následujících fázích.

Typy tokenů jsou definovány výčtem \texttt{TokenType} v~\texttt{src/types.rs}.
Lexer rozpoznává celočíselné literály (\texttt{i32}), identifikátory (slova),
oddělovače definic (\texttt{:} a~\texttt{;}), komentáře v~závorkách
\texttt{(\dots)} a~řetězcové literály.

Specifickým prvkem inspirovaným \Forth{} je podpora syntaxe \texttt{S\" ...\"},
která umožňuje zapisovat řetězec bez explicitního uzavírání do tokenů jednotlivých
slov. Lexer tento vzor detekuje ještě před obecným rozlišením tokenu.

\begin{lstlisting}[language=Rust, caption={Rozpoznání tokenu v lexeru}, float=htbp]
fn next_token(&mut self, start_pos: Position) -> Result<Token, ParseError> {
    let ch = self.current_char();
    if (ch == 'S' || ch == 's') && self.peek_char() == Some('"') {
        return self.read_s_quote_literal(start_pos);
    }
    match ch {
        '(' => self.read_comment(start_pos),
        '"' => self.read_string_literal(start_pos),
        ':' | ';' => self.read_definition_token(start_pos),
        _ => self.read_token(start_pos),
    }
}
\end{lstlisting}

Metoda \texttt{read\_token} načte nejdelší podřetězec oddělený bílými znaky
a~pokusí se jej převést na číslo. Pokud převod selže, token je interpretován jako
slovo a~normalizován na velká písmena (\texttt{to\_uppercase}). To sjednocuje
zpracování slov nezávisle na velikosti písmen ve zdrojovém kódu.

Komentáře v~závorkách jsou tokenizovány jako \texttt{TokenType::Comment}.
Parser i~další fáze je následně mohou ignorovat nebo využít (např. pro zachování
informativních poznámek v~IR ve formě \texttt{Comment} instrukcí).
